\documentclass[instructions.tex]{subfiles}
\begin{document}
 
\chapter{Aveuglement de ceux qui s'oublient eux-mêmes pour enrichir leurs héritiers.}
 
\primo Quelle folie d'amasser du bien sans savoir pour qui! Vous épargnez pour vos enfans ou pour vos parens; mais peut-être, deux mois après votre mort, vos biens passeront en d'autres mains, peut-être à vos ennemis, qui jouiront de ces biens pendant que vous gémirez dans les flammes, pour les avoir conservés ou amassés avec trop d'attache. Quand vous seriez assuré de vos héritiers, êtes-vous assuré qu'ils useront sagement de ce que vous aurez laissé? La débauche, les procès, le jeu, ne dévoreront-ils pas le fruit de vos épargnes? Ce que vous aurez amassé par le péché, ne se dissipera-t-il point par d'autres péchés, qui attireront sur vos héritiers et sur vous la colère de Dieu?
 
Oh! qu'on est aveugle, lorsque, par une amitié trop humaine, on est plus sensible aux commodités d'autrui qu'au salut de son âme, et qu'on donne aux autres ce qu'on refuse à soi-même, sans en recevoir ni récompense ni retour! Ne traiteriez-vous pas d'insensé, dit Salvien, un homme qui raisonnerait ainsi : Mes chers héritiers, je ne veux pas me servir de mes biens ni gagner le ciel par mes aumônes; je veux renoncer aux satisfactions de la vie, et m'exposer à être damné, afin que vous soyez plus à votre aise? Peut-on penser et agir avec plus d'extravagance? Ce n'est plus une amitié, mais une fureur, que de se perdre pour faire plaisir aux autres.
 
Il est bien plus sûr de faire part de ses biens à Jésus-Christ en la personne des pauvres, pour en recevoir la récompense dans le ciel, que de les laisser à des héritiers qui ne peuvent nous rendre heureux, quand même ils le voudraient.
 
\secundo On entend tous les jours des pères de famille faire cette demande : Ne ferai-je pas tort à mes enfans, de ne pas leur conserver mon bien? Il y a de l'équivoque dans cette demande : elle est souvent le prétexte d'un père avare et chicaneur. Oui, vous feriez tort à vos enfans de dissiper vos biens en jeux, en débauches, en procès. Vous feriez tort à vos enfans de leur laisser des biens injustement acquis, mais vous ne leur faites aucun tort en faisant des aumônes, ni en cédant à un plaideur obstiné. Dieu vous a donné des enfans, non pas précisément pour leur laisser du bien, mais pour les sanctifier.
 
Un père, à la vérité, doit user d'une sage économie pour établir ses enfans; mais s'il est coupable lorsqu'il abandonne l'éducation et le soin de sa famille, il n'est pas excusable lorsqu'il abandonne le pauvre. N'est-ce pas une criante injustice, dit saint Augustin, de laisser languir Jésus-Christ dans le pauvre, tandis que rien ne manque aux débauches de votre fils\footnote{Injustitia magna est ut egeat Dominus tuus, et habeat unde luxuriétur filius tuus. De Discip. 3.}? Vous êtes en peine de ce que deviendront vos enfans après votre mort : soyez bien plus en peine de ce que vous deviendrez vous-même. Sanctifiez-vous, et tâchez de mériter la miséricorde de Dieu, en prenant soin de vos enfans. Le moyen le plus sûr de les enrichir, dit saint Chrysostome, c'est de faire l'aumône.
 
Vos richesses, mises dans le sein des pauvres, sont en assurance pour vous, et loin d'être perdues pour vos enfans, elles sont, dit saint Eucher, un trésor mieux placé qu'entre leurs mains\footnote{Thesaurum suum is benè recondit, qui indigentibus dividit. De ver. Sap.}. Sainte Marcelle, illustre dame romaine, dépouillait ses enfans pour revêtir et nourrir les pauvres : Ne craignez rien, mes enfans, leur disait-elle, je vous laisse dans la miséricorde de Jésus-Christ un fonds plus assuré que tous nos biens. Qu'on se souvienne qu'une famille fondée sur l'aumône ne périra pas.
 
Après tout, il s'agit de vous sauver. Si vous n'avez pas soin de vous-même pendant votre vie, qui en aura soin après votre mort? Vos enfans pourront-ils vous tirer de l'enfer? Hélas! s'écrie Salvien, toutes les richesses dans les mains des héritiers ne sont pas capables d'éteindre une étincelle des flammes qui brûlent ces morts infortunés\footnote{Flammae infelicium mortuorum divitiis non refrigerantur haeredum. Lib. 8 ad Eccl. cat.}.
 
Un service, quelques messes, voilà tout ce que vous pouvez attendre de vos héritiers. Il vaut mieux qu'il leur manque quelques biens, et qu'il ne manque rien pour votre salut. On leur laisse toujours assez, quand ce qu'on leur donne est précédé des exemples de vertu, et qu'on leur laisse la crainte de Dieu. N'estimez les richesses qu'autant qu'elles servent pour gagner le ciel; vous éprouverez que l'unique consolation du riche, c'est d'être libéral et charitable.
 
\section{Résolutions}
 
\primo J'userai de ce que Dieu m'a donné, selon l'ordre de son adorable providence. 

\secundo Je me donnerai à moi-même ce que la nécessité exige.

\tertio Je me ferai des amis dans le ciel par mes saintes largesses envers les pauvres. 

\quarto Et je pourvoirai à mon salut avant que de penser à enrichir mes héritiers. 
\prayer{Faites, ô mon Dieu, que je ne me perde pas par le mauvais usage des biens périssables de ce monde.}
 
\end{document}
